\chapter{Teoretická část studentské práce}

%Teoretické zázemí studentské práce vhodně rozdělené do částí.

%(Struktura navržená v~této šabloně je nejhrubší možná, po konzultaci s~vedoucím je vhodné zvolit přiléhavější.)

%% Dopravnikovy pas
\section{Pásový dopravník}
Přeprava materiáu nebo předmětů po dopravních pásech je jeden z nejrozšířenějších a nejefektivnějších způsobů souběžné (kontinuální) přepravy. Poskytují plynulý přesun výrobků mezi jednotlivými pracovišti. 
[doplnit citaci ze je to nejrozšířenější]
V rámci návrhu fyzického modelu pro tuto semestrální práci bylo nutné zajistit funkční a řiditelný dopravní mechanismus. V následujících podkapitolách je podrobnější analýza komerčně dostupných průmyslových řešení a následně je popsána vlastní realizace, která z důvodu udržení nízkých finančních nákladů spočívá v modifikaci fyzického hardwaru funkčního domácího běžeckého pásu. 
\subsection{Průzkum trhu a dostupná řešení}

\subsection{Výběr a modifikace běžeckého pásu}

\subsection{Bezdrátová komunikace v IoT}

\cite{IR}
\cite{IoT}

%% Kamera

\section{Smart kamera}

\subsection{Podporované funkce kamery}

\subsection{aa}

%% Automation Studio

\section{Prostředí Automation Studio}

\subsection{Popis projektu s Deltabotem}